% fig_vl05_diagonal.tex — Diagonalisierung als Basiswechsel: derselbe
% Operator, einmal in der Standardbasis, einmal in der Eigenbasis.
%
% Follows the conventions of fig_vl11_kasten.tex (the reference plate).
%
%   * geometry comes from the lecture's own worked example and nothing
%     else.  Vorlesung 5 diagonalises A = sigma_x = [[0,1],[1,0]] with the
%     Hadamard matrix U = (1/sqrt2)[[1,1],[1,-1]].  Drawn as it stands,
%     sigma_x is a reflection: it maps the unit circle onto itself and
%     there is no shape to see.  So the plate draws the shifted operator
%         Atilde = A + 2*Id = [[2,1],[1,2]],
%     which has exactly the same eigenvectors and is diagonalised by
%     exactly the same U from the lecture's Beispiel, its eigenvalues
%     merely raised by 2:
%         |phi_1> = (1,1)/sqrt2,  a_1 =  1 + 2 = 3
%         |phi_2> = (1,-1)/sqrt2, a_2 = -1 + 2 = 1
%         U Atilde U^dag = diag(3,1).
%     Every label on the plate says \tilde A and never A, because the
%     lecture's A is sigma_x and stands three lines above the figure; the
%     prose that introduces the figure defines the shift.
%     Every curve below is that matrix applied to the unit circle,
%     plotted parametrically, not sketched:
%         panel (a)  t |-> (2cos t + sin t,  cos t + 2 sin t)
%         panel (b)  t |-> (3cos t,          sin t)
%     and the two are the same ellipse, seen from the two bases: U carries
%     (a) into (b).  (U is a reflection, det U = -1, not a rotation.  On
%     the centrally symmetric ellipse a -45 degree turn would give the
%     same curve, but not the same arrows: it sends |phi_2> to -e_2, where
%     U sends it to +e_2, and +e_2 is what is drawn.)  Semi-axes 3 and 1
%     are the eigenvalues, principal axes are the eigenvectors; both
%     follow from Atilde = Atilde^dag.
%     The arrow tips are the images of the basis vectors and therefore the
%     columns of the matrix: Atilde|1> = (2,1), Atilde|2> = (1,2) in (a),
%     Atilde|phi_1> = 3|phi_1>, Atilde|phi_2> = |phi_2> in (b).
%
%   * both panels are drawn at the same units-per-cm and on the same
%     ranges, because the whole claim of the figure is that the object is
%     the same and only the frame moved.  A different scale in the two
%     panels would destroy exactly that reading.
%
%   * colour encodes *which basis vector*, and nothing else: cgblue for
%     the first, birrenorange for the second.  Input and image are told
%     apart by tint and weight rather than by hue — the incoming unit
%     vector is a wide pale arrow, its image a narrower saturated one, so
%     the pair survives greyscale printing and, where the two coincide
%     (a_2 = 1 in panel (b)), the pale halo shows the coincidence instead
%     of hiding it.  The unit circle, the coordinate rays and the
%     principal axes are neutral grey: they are reference lines, not
%     data.  The image ellipse is neutral navy — it belongs to no basis,
%     which is the point.
%
%   * the dotted tip-to-tip segments in (a) say which point of the circle
%     goes to which point of the ellipse.  In (b) they are omitted: there
%     the displacement runs along the image arrow itself and a dotted
%     line on top of it would only be noise.
%
% Panel (a) standard basis.  The matrix has off-diagonal entries, the
%           ellipse lies askew, and neither basis vector comes back
%           parallel to itself.
% Panel (b) eigenbasis.  Same ellipse, axes now on the coordinate axes,
%           matrix diagonal, and the arrows only stretch: |phi_1> by the
%           factor 3, |phi_2> not at all.  Circle and ellipse touch at
%           +/- |phi_2> precisely because a_2 = 1.

\begin{center}
\begin{tikzpicture}[
    refray/.style   ={-{Stealth[length=4pt,width=3pt]}, line width=0.5pt,
                      birrengray!60},
    refcurve/.style ={line width=0.7pt, birrengray!70, densely dashed},
    princ/.style    ={line width=0.6pt, birrengray!70,
                      dash pattern=on 3.5pt off 1.5pt on 0.8pt off 1.5pt},
    imgcurve/.style ={line width=1.2pt, spacecadet},
    invec/.style    ={-{Stealth[length=6pt,width=5.5pt]}, line width=2.4pt},
    outvec/.style   ={-{Stealth[length=5.5pt,width=4pt]}, line width=1.3pt},
    movedot/.style  ={line width=0.7pt, densely dotted},
    % The four ket labels sit inside the unit circle, where the circle and
    % the principal axes leave no gap wide enough for a two-glyph label.
    % They therefore carry the same faint white backing the house legend
    % style uses, so a reference line never runs through a bracket.
    ketlbl/.style   ={lbl, font=\scriptsize, inner sep=0.9pt,
                      fill=white, fill opacity=0.85, text opacity=1},
]

% ─────────────────────────────────────────────────────────────
%  (a)  Standard basis: Atilde = [[2,1],[1,2]] is not diagonal
% ─────────────────────────────────────────────────────────────
\begin{axis}[
    name=std,
    tpcplot,
    scale only axis,
    width=5.7cm, height=5.7cm,     % equal spans below: 45 degrees looks like 45
    axis lines=none, grid=none,
    xmin=-3.35, xmax=3.35,
    ymin=-3.35, ymax=3.35,
    clip=false,
    title={(a) \tpcifenglish{Standard basis}{Standardbasis}},
]

  % ---- coordinate frame (reference) ----
  \draw[refray] (axis cs:-2.60,0) -- (axis cs:3.15,0);
  \draw[refray] (axis cs:0,-2.60) -- (axis cs:0,3.15);

  % ---- unit circle (reference): the input ----
  \addplot[refcurve, variable=t, domain=0:360, samples=181]
    ({cos(t)},{sin(t)});

  % ---- principal axes of the image ellipse = the eigen-directions,
  %      drawn only over the ellipse: semi-axes a_1 = 3, a_2 = 1 ----
  \draw[princ] (axis cs:-2.1213,-2.1213) -- (axis cs:2.1213,2.1213);
  \draw[princ] (axis cs:-0.7071,0.7071)  -- (axis cs:0.7071,-0.7071);
  \node[lbl, text=birrengray, font=\scriptsize, anchor=north east]
    at (axis cs:-2.16,-2.16) {$|\phi_1\rangle$};
  \node[lbl, text=birrengray, font=\scriptsize, anchor=north west]
    at (axis cs:0.74,-0.78) {$|\phi_2\rangle$};

  % ---- image of the unit circle under A ----
  \addplot[imgcurve, variable=t, domain=0:360, samples=181]
    ({2*cos(t)+sin(t)},{cos(t)+2*sin(t)});

  % ---- where the tips of the basis vectors go ----
  \draw[movedot, cgblue]       (axis cs:1,0) -- (axis cs:2,1);
  \draw[movedot, birrenorange] (axis cs:0,1) -- (axis cs:1,2);

  % ---- the basis vectors and their images ----
  \draw[invec,  cgblue!30]       (axis cs:0,0) -- (axis cs:1,0);
  \draw[invec,  birrenorange!30] (axis cs:0,0) -- (axis cs:0,1);
  \draw[outvec, cgblue]          (axis cs:0,0) -- (axis cs:2,1);
  \draw[outvec, birrenorange]    (axis cs:0,0) -- (axis cs:1,2);

  % |1> sits outside the ellipse to the lower right of its own arrowhead:
  % every spot nearer the origin is crossed by a ray or by the ellipse.
  \node[ketlbl, text=cgblue, anchor=north west]
    at (axis cs:1.32,-0.08) {$|1\rangle$};
  \node[ketlbl, text=birrenorange, anchor=east]
    at (axis cs:-0.12,0.70) {$|2\rangle$};
  \node[lbl, text=cgblue, font=\scriptsize, anchor=west]
    at (axis cs:2.10,0.92) {$\tilde A|1\rangle$};
  \node[lbl, text=birrenorange, font=\scriptsize, anchor=south]
    at (axis cs:1.00,2.14) {$\tilde A|2\rangle$};

  % ---- the matrix this frame produces ----
  % Two lines: on one line the definition and the matrix together reach
  % as far as the note in the opposite corner.
  \node[lbl, anchor=north west, align=left]
    at (axis cs:-3.30,3.22)
    {$\tilde A = A + 2\mathbb{I}$\\[3pt]
     $\phantom{\tilde A} = \begin{pmatrix} 2 & 1 \\ 1 & 2\end{pmatrix}$};

  \node[lbl, text=birrengray, font=\scriptsize, anchor=north east]
    at (axis cs:3.30,3.18)
    {\tpcifenglish{circle $\mapsto$ ellipse}{Kreis $\mapsto$ Ellipse}};

  \node[lbl, text=birrengray, font=\scriptsize, anchor=north]
    at (axis cs:0,-2.74)
    {\tpcifenglish{no basis vector is an eigenvector}%
                  {kein Basisvektor ist Eigenvektor}};
\end{axis}

% ─────────────────────────────────────────────────────────────
%  (b)  Eigenbasis: same ellipse, U Atilde U^dag = diag(3,1)
% ─────────────────────────────────────────────────────────────
\begin{axis}[
    tpcplot,
    scale only axis,
    at={(std.south east)}, anchor=south west, xshift=20mm,
    width=5.7cm, height=5.7cm,
    axis lines=none, grid=none,
    xmin=-3.35, xmax=3.35,
    ymin=-3.35, ymax=3.35,
    clip=false,
    title={(b) \tpcifenglish{Eigenbasis}{Eigenbasis}},
]

  % ---- coordinate frame: now the eigen-directions ----
  \draw[refray] (axis cs:-2.60,0) -- (axis cs:3.25,0);
  \draw[refray] (axis cs:0,-2.60) -- (axis cs:0,3.15);

  % ---- unit circle (reference) ----
  \addplot[refcurve, variable=t, domain=0:360, samples=181]
    ({cos(t)},{sin(t)});

  % ---- the same ellipse, now axis-aligned: semi-axes a_1 = 3, a_2 = 1 ----
  \addplot[imgcurve, variable=t, domain=0:360, samples=181]
    ({3*cos(t)},{sin(t)});

  % ---- the basis vectors and their images: pure stretching ----
  \draw[invec,  cgblue!30]       (axis cs:0,0) -- (axis cs:1,0);
  \draw[invec,  birrenorange!30] (axis cs:0,0) -- (axis cs:0,1);
  \draw[outvec, cgblue]          (axis cs:0,0) -- (axis cs:3,0);
  \draw[outvec, birrenorange]    (axis cs:0,0) -- (axis cs:0,1);

  % Both labels are kept inside the unit circle: at their old places the
  % dashed circle ran straight through the ket brackets.
  \node[ketlbl, text=cgblue, anchor=north]
    at (axis cs:0.45,-0.14) {$|\phi_1\rangle$};
  \node[ketlbl, text=birrenorange, anchor=east]
    at (axis cs:-0.10,0.58) {$|\phi_2\rangle$};

  \node[lbl, text=cgblue, font=\scriptsize, anchor=north]
    at (axis cs:1.75,-1.10) {$\tilde A|\phi_1\rangle = 3\,|\phi_1\rangle$};
  \node[lbl, text=birrenorange, font=\scriptsize, anchor=south west]
    at (axis cs:0.14,1.10) {$\tilde A|\phi_2\rangle = |\phi_2\rangle$};

  % ---- the matrix this frame produces ----
  \node[lbl, anchor=north west] at (axis cs:-3.30,3.22)
    {$U\tilde AU^\dagger
       = \begin{pmatrix} 3 & 0 \\ 0 & 1\end{pmatrix}$};

  \node[lbl, text=birrengray, font=\scriptsize, anchor=north]
    at (axis cs:0,-2.74)
    {\tpcifenglish{every basis vector is an eigenvector}%
                  {jeder Basisvektor ist Eigenvektor}};
\end{axis}

% ─────────────────────────────────────────────────────────────
%  The basis change between the two frames.  U is the Hadamard matrix
%  the lecture's Beispiel already computed.
% ─────────────────────────────────────────────────────────────
\draw[-{Stealth[length=5pt,width=3.5pt]}, spacecadet, line width=0.8pt]
  ([xshift=6mm,yshift=8mm]std.east) -- ([xshift=14mm,yshift=8mm]std.east);
\node[lbl, font=\scriptsize, anchor=south]
  at ([xshift=10mm,yshift=8.6mm]std.east) {$U$};

\end{tikzpicture}
\end{center}
